% Minimal retained preamble for lane-p2 cards taken from
% sources/logic-fl05-ramsey-like-theorems-for-the-schre/pinned-original-source.tex
% (arXiv:2412.11598v3, CC BY 4.0).  This source ships no separate preamble file:
% the pinned tex is a concatenation of the authors' main.tex and macros.tex and
% carries its preamble inline, so the blocks below are sliced mechanically from
% that file and nothing has been renamed or redefined.
%
% Copied verbatim from the source:
%   line 2        : the authors' class line
%   lines 3-5,7,9 : graphicx, amsthm, the ams bundle, cleveref, url
%   line 10       : tikz-cd, used by the authors' commutative diagrams
%   lines 20-108  : the authors' macro definitions, taken in full so that every
%                   macro of the retained spans is available
%   lines 113-124 : the authors' theorem declarations
%   lines 126-128 : the authors' macros for the three calligraphic letters, which
%                   sit just after the theorem declarations in the source
%   lines 131-141 : the authors' machinery for restating a theorem
%   lines 143-148 : the authors' coloured comment macros
%
% Omitted source preamble lines and why:
%   line 1, 110-111 and 150 (concatenation comment, title, date, author): the
%       assembler supplies the title block, author list and date;
%   line 6 (hyperref) and line 8 (inputenc with the utf8x option): the forced
%       xelatex engine supplies the hyperref wiring itself and utf8x is
%       incompatible with it;
%   line 115 (numberwithin theorem by section): the assembler numbers the
%       retained material itself;
%   lines 120 and 121-123 are declarations that the authors left commented out or
%       that duplicate earlier ones and are not reproduced.
\documentclass{article}
\usepackage{graphicx}
\usepackage{amsthm}
\usepackage{amsfonts, amsmath, amssymb, amsthm}
\usepackage{cleveref}
\usepackage{url}
\usepackage{tikz-cd}
\newcommand{\Psf}{\mathsf{P}}
\newcommand{\Qsf}{\mathsf{Q}}
\newcommand{\Rsf}{\mathsf{R}}
\def\F{\mathcal{F}}
\def\U{\mathcal{U}}
\def\C{\mathcal{C}}
\def\I{\mathcal{I}}
\def\D{\mathcal{D}}
\def\L{\mathcal{L}}
\def\V{\mathcal{V}}
\def\A{\mathcal{A}}
\def\B{\mathcal{B}}
\def\P{\mathcal{P}}
\def\M{\mathcal{M}}
\def\Nc{\mathcal{N}}
\def\Q{\mathcal{Q}}
\newcommand{\cs}{2^{\NN}}
\newcommand{\bs}{\NN^{\NN}}
\newcommand{\bstr}{2^{<\NN}}
\newcommand{\uh}{\upharpoonright}
\newcommand{\converge}{\!\!\downarrow}
\newcommand{\diverge}{\!\!\uparrow}
\newcommand{\qvdash}{\operatorname{{?}{\vdash}}}
\newcommand{\nqvdash}{\operatorname{{?}{\nvdash}}}
\newcommand{\ISig}{\mathsf{I}\Sigma^0}
\newcommand{\Robinson}{\mathsf{Q}}
\newcommand{\IPi}{\mathsf{I}\Pi^0}
\newcommand{\IDelta}{\mathsf{I}\Delta^0}
\newcommand{\Iopen}{\mathsf{I_{open}}}
\newcommand{\BSig}{\mathsf{B}\Sigma^0}
\newcommand{\BPi}{\mathsf{B}\Pi^0}
\newcommand{\BDelta}{\mathsf{B}\Delta^0}
\newcommand{\BCSig}{\mathsf{BC}\Sigma^0}
\newcommand{\BCPi}{\mathsf{BC}\Pi^0}
\newcommand{\BCDelta}{\mathsf{BC}\Delta^0}
\newcommand{\LSig}{\mathsf{L}\Sigma^0}
\newcommand{\LPi}{\mathsf{L}\Pi^0}
\newcommand{\CSig}{\mathsf{C}\Sigma}
\newcommand{\RCA}[0]{\mathsf{RCA}}
\newcommand{\WKL}[0]{\mathsf{WKL}}
\newcommand{\WWKL}[0]{\mathsf{WWKL}}
\newcommand{\ACA}[0]{\mathsf{ACA}}
\newcommand{\ATR}[0]{\mathsf{ATR}}
\newcommand{\PIOOCA}[0]{\Pi^1_1\mathsf{-CA}}
\newcommand{\CAC}[0]{\mathsf{CAC}}
\newcommand{\ADS}[0]{\mathsf{ADS}}
\newcommand{\RRT}[0]{\mathsf{RRT}}
\newcommand{\FS}[0]{\mathsf{FS}}
\newcommand{\PFS}[0]{\mathsf{PFS}}
\newcommand{\TS}[0]{\mathsf{TS}}
\newcommand{\GRRT}[0]{\mathsf{GRRT}}
\newcommand{\LRRT}[0]{\mathsf{LRRT}}
\newcommand{\EM}[0]{\mathsf{EM}}
\newcommand{\SEM}[0]{\mathsf{SEM}}
\newcommand{\HalfEM}[0]{\mathsf{Half}\mbox{-}\mathsf{EM}}
\newcommand{\RT}[0]{\mathsf{RT}}
\newcommand{\SRT}[0]{\mathsf{SRT}}
\newcommand{\NN}[0]{\mathbb{N}}
\newcommand{\nat}[0]{\mathbb{N}}
\newcommand{\card}{\operatorname{card}}
\newcommand{\finsub}{\subseteq_{\mathtt{fin}}}
\newcommand{\W}{\mathrm{W}}
\newcommand{\sW}{\mathrm{sW}}
\newcommand{\lex}{\mathrm{lex}}
\newcommand{\base}{\operatorname{base}}
\newcommand{\range}{\operatorname{ran}}
\def\qt#1{``#1''}%

\newtheorem*{statement}{Statement}
\newtheorem{theorem}{Theorem}
\newtheorem{maintheorem}[theorem]{Main Theorem}
\newtheorem{lemma}[theorem]{Lemma}
\newtheorem{question}[theorem]{Question}
\newtheorem{proposition}[theorem]{Proposition}
%\newtheorem{statement}[theorem]{Statement}
\newtheorem{remark}[theorem]{Remark}
\newtheorem{definition}[theorem]{Definition}
\newtheorem{corollary}[theorem]{Corollary}
\newtheorem{example}[theorem]{Example}
\newcommand{\R}{\mathcal{R}}
\renewcommand{\S}{\mathcal{S}}
\newcommand{\T}{\mathcal{T}}

% Allow to repeat theorems
\makeatletter
\newtheorem*{rep@theorem}{\rep@title}
\newcommand{\newreptheorem}[2]{%
\newenvironment{rep#1}[1]{%
 \def\rep@title{#2 \ref{##1}}%
 \begin{rep@theorem}}%
 {\end{rep@theorem}}}
\makeatother

\newreptheorem{theorem}{Theorem}

% Comments
\usepackage{xcolor}
\newcommand{\ludovic}[1]{\textcolor{blue}{#1}}
\newcommand{\quentin}[1]{\textcolor{green}{#1}}
\newcommand{\lorenzo}[1]{\textcolor{orange}{#1}}
\newcommand{\oriola}[1]{\textcolor{red}{#1}}
\crefname{originallemma}{original lemma}{original lemmas}
\Crefname{originallemma}{Original lemma}{Original lemmas}
\makeatletter
\newlabel{lem:fs-forcing-question-sigma01}{{6.23}{source}{}{sourcepointer-1}{}}
\newlabel{lem:fs-forcing-question-sigma01@cref}{{[originallemma][0][]6.23}{[1][0][]source}{}{}{}}
\makeatother
\crefname{originallemma}{original lemma}{original lemmas}
\Crefname{originallemma}{Original lemma}{Original lemmas}
\makeatletter
\newlabel{lem:fs-stabilizing}{{6.18}{source}{}{sourcepointer-2}{}}
\newlabel{lem:fs-stabilizing@cref}{{[originallemma][0][]6.18}{[1][0][]source}{}{}{}}
\makeatother
\crefname{originallemma}{original lemma}{original lemmas}
\Crefname{originallemma}{Original lemma}{Original lemmas}
\makeatletter
\newlabel{lem:fs-preserving-b}{{6.20}{source}{}{sourcepointer-3}{}}
\newlabel{lem:fs-preserving-b@cref}{{[originallemma][0][]6.20}{[1][0][]source}{}{}{}}
\makeatother
\crefname{originallemma}{original lemma}{original lemmas}
\Crefname{originallemma}{Original lemma}{Original lemmas}
\makeatletter
\newlabel{lem:fs-preserving-a}{{6.16}{source}{}{sourcepointer-4}{}}
\newlabel{lem:fs-preserving-a@cref}{{[originallemma][0][]6.16}{[1][0][]source}{}{}{}}
\makeatother
\title{The forcing question meets its specification: it decides which extensions force a sigma-0-1 formula and its negation}
\author{Lorenzo Carlucci, Oriola Gjetaj, Quentin Le Houérou, Ludovic Levy Patey}
\date{arxiv:2412.11598v3; CC BY 4.0}
\begin{document}
\maketitle
\noindent\textit{Editorial note.} The counted unit is the article's own Lemma lem:fs-forcing-question-spec at source lines 1221-1227 together with its proof at source lines 1228-1236, a straight adjacent pairing in the pinned file. The pinned file is a concatenation of the authors' main and macro files; it is the article's fixed version arXiv:2412.11598v3 under CC BY 4.0.\par\medskip
\noindent\textit{Editorial note.}  The definition of the forcing question is retained complete, from its opening marker at line 1196 through its closing marker at line 1199; the first revision cut the span one line short and the assembler rightly refused to close the environment, so this revision carries the corrected span.\par\medskip
\noindent\textit{Editorial note.} The retained preamble is the lane's shared slice for this source, now 137 lines: it carries the authors' macro block together with the three calligraphic letter macros that sit between the macro block and the theorem declarations, which this lemma's statement needs for the space of relations.\par\medskip
\noindent\textit{Editorial note.} Two context spans carry the notions the retained statement speaks about: the definition of the forcing question at lines 1196-1199 and the definition of forcing for a condition at line 1218. Both are context, not counted units, and neither overlaps the counted interval. They are deliberately kept out of the interval handed to the gate, because that interval is meant to be the counted unit itself and the gate's label guard refuses an interval that swallows other labelled environments.\par\medskip
\noindent\textit{Editorial note.} The retained proof names four lemmas of the same section at source lines 1121-1125, 1144-1147, 1162-1167 and 1201-1204, each counted by a separate card of this lane. All four are declared as non-load-bearing navigation pointers with the article's own section-wise numbers 6.16, 6.18, 6.20 and 6.23, counted over the numbered environments of section 6; none of their proofs is reproduced and no bibliography is needed because the retained spans contain no citation.\par\medskip
\noindent\textit{Editorial note.} The second direction of the proof introduces a thinned colouring and checks four properties of it in one paragraph, then argues by contradiction through the defining property (b); all of that is the author's own argument and is retained line for line.\par\medskip
\noindent\textit{Allowed standard prerequisites (editorial).}\begin{itemize}
\item Conditions, the forcing question, forcing and the effectiveness of the forcing question. A condition is a triple of a colouring constrained by a computable bounding function, a finite stem and an infinite reservoir computing a fixed set D, with the defining properties (a) and (b); the forcing question holds for a condition and a formula when every relation of the bounded space admits a finite free block in the reservoir that makes the formula true on the enlarged stem, and the lemma immediately preceding this one shows that question to be equivalent to a first-order statement about a single level of bounded progressive functions. Forcing for a condition is the property that every sufficiently generic filter containing it makes the formula true of the induced set. The author's own definitions and that preceding lemma are the allowed inputs.. The counted lemma is the specification of the forcing question: it is what licenses reading the first-order question as a statement about extensions that force the formula or its negation, and it is invoked by the diagonalisation that follows.
\end{itemize}\medskip
\ifdefined\hypertarget\hypertarget{sourcepointer-1}{}\fi
\noindent\textit{Source reference.} Original lemma 6.23: sources/logic-fl05-ramsey-like-theorems-for-the-schre/pinned-original-source.tex, source lines 1201--1204. Navigation pointer invoked by the retained proof to turn the forcing question into a first-order statement about one level. That lemma is counted by a separate card of this lane, so its proof is not reproduced here. The number is the article's own section-wise number for it, counted over the numbered environments of section 6.\par\medskip
\ifdefined\hypertarget\hypertarget{sourcepointer-2}{}\fi
\noindent\textit{Source reference.} Original lemma 6.18: sources/logic-fl05-ramsey-like-theorems-for-the-schre/pinned-original-source.tex, source lines 1144--1147. Navigation pointer invoked by the retained proof for the stabilising reservoir at the level supplied by the previous pointer. That lemma is counted by a separate card of this lane. The number is the article's own section-wise number, counted over the numbered environments of section 6.\par\medskip
\ifdefined\hypertarget\hypertarget{sourcepointer-3}{}\fi
\noindent\textit{Source reference.} Original lemma 6.20: sources/logic-fl05-ramsey-like-theorems-for-the-schre/pinned-original-source.tex, source lines 1162--1167. Navigation pointer invoked by the retained proof to keep property (b) when the witnessing block is appended in the first direction, and again implicitly in the second. That lemma is counted by a separate card of this lane. The number is the article's own section-wise number, counted over the numbered environments of section 6.\par\medskip
\ifdefined\hypertarget\hypertarget{sourcepointer-4}{}\fi
\noindent\textit{Source reference.} Original lemma 6.16: sources/logic-fl05-ramsey-like-theorems-for-the-schre/pinned-original-source.tex, source lines 1121--1125. Navigation pointer invoked by the retained proof in both directions to restore property (a) by shrinking the reservoir. That lemma is counted by a separate card of this lane. The number is the article's own section-wise number, counted over the numbered environments of section 6.\par\medskip
\begin{definition}
Let $(f, \sigma, X)$ be a condition and $\varphi(G)$ be a $\Sigma^0_1$-formula.
Let $(f, \sigma, X) \qvdash \varphi(G)$ hold if and only if for every relation $R \in \R_{b_f^{+}}$, there is a finite $g_R$-free set $\rho \subseteq X$ such that $\varphi(\sigma \cup \rho)$ holds.
\end{definition}

We say that a condition $p$ \emph{forces} a formula $\varphi(G)$ if $\varphi(G_\F)$ holds for every sufficiently generic filter~$\F$ containing~$p$.

\begin{lemma}\label[lemma]{lem:fs-forcing-question-spec}
Let $p = (f, \sigma, X)$ be a condition and $\varphi(G)$ be a $\Sigma^0_1$-formula.
\begin{enumerate}
    \item[1.] If $p \qvdash \varphi(G)$, then there is an extension of~$p$ forcing $\varphi(G)$.
    \item[2.] If $p \nqvdash \varphi(G)$, then there is an extension of~$p$ forcing $\neg \varphi(G)$.
\end{enumerate}
\end{lemma}

\begin{proof}
Suppose first $p \qvdash \varphi(G)$ holds. By \Cref{lem:fs-forcing-question-sigma01}, there is some~$k \in \NN$ such that for every $b_f^{+}$-constrained progressive function $g : [0, k]^{\leq h(\cdot)} \to [0, k]^{<\omega}$, there is a finite $g$-free set $\rho \subseteq X$ such that $\varphi(\sigma \cup \rho)$ holds. By \Cref{lem:fs-stabilizing}, there is an infinite subset $Y_0 \subseteq X$ stabilizing $[0, k]$ and such that $D \not \leq_T Y_0$. Let $g_{k,Y_0} : [0, k]^{\leq h(\cdot)} \to [0, k]^{<\omega}$ be the limit function. Note that $g_{k,Y_0}$ is $b_f^{+}$-constrained and progressive, so there is a finite $g_{k,Y_0}$-free set $\rho \subseteq X$ such that $\varphi(\sigma \cup \rho)$ holds.
By \Cref{lem:fs-preserving-b}, $(f, \sigma \cup \rho, Y_0)$ satisfies (b).
By \Cref{lem:fs-preserving-a}, there is an infinite subset $Y \subseteq Y_0$ such that $(f, \sigma \cup \rho, Y)$ satisfies (a) and $D \not \leq_T Y$. Thus $(f, \sigma \cup \rho, Y)$ is a valid extension of~$p$. By choice of~$\rho$, it forces $\varphi(G)$.

Suppose $p \nqvdash \varphi(G)$ holds. Then there is a relation $R \in \R_{b_f^{+}}$ such that for every finite $g_R$-free set $\rho \subseteq X$, $\varphi(\sigma \cup \rho)$ does not hold. Let $g_R \in \P_{b_f^{+}}$ be the corresponding function. Let $\hat{f} : [\NN]^{\leq h(\cdot)} \to [\NN]^{<\omega}$ be defined by $\hat{f}(s) = f(s) \cup (g_R(s) \setminus \sigma)$. Note that $\hat{f} \leq f$, and every $\hat{f}$-free subset $\rho \subseteq X$ is $g_R$-free. Moreover, $\hat{f}$ is $(b_f + b_f^{+})$-constrained and progressive. Since $(f, \sigma, X)$ satisfies (b) and for every $s \in [\sigma \cup X]^{\leq h(\cdot)}$, $\hat{f}(s)\cap \sigma\subseteq f(s)\cap \sigma\subseteq s$, then $(\hat{f}, \sigma, X)$ satisfies (b). By \Cref{lem:fs-preserving-a}, there is an infinite subset $Y \subseteq X$ such that $D \not \leq_T Y$ and $(\hat{f}, \sigma, Y)$ satisfies (a). Thus $q = (\hat{f}, \sigma, Y)$ is a valid extension of~$p$. 

We claim that $q$ forces $\neg \varphi(G)$. Indeed, suppose there is an extension $(\tilde{f}, \tau, Z) \leq q$ such that $\varphi(\tau)$ holds. Then by property (b) of  $(\tilde{f}, \tau, Z)$, $\tau$ is $\tilde{f}$-free, and since $\tilde{f} \leq \hat{f}$, $\tau$ is $\hat{f}$-free. Let $\rho = \tau \setminus \sigma$. By definition of~$\hat{f}$, $\rho$ is a $g_R$-free subset of~$X$, contradicting our choice of~$g_R$.
\end{proof}

\end{document}
